\section{Introduction}\label{sec:introduction}

A discrete group $G$ is \emph{amenable} if, for every finite
$K\subseteq G$ and every $\varepsilon>0$, there is a nonempty finite
$D\subseteq G$ such that
\begin{equation}\label{eq:folner-introduction}
 |gD\mathbin{\triangle}D|<\varepsilon |D|
 \qquad(g\in K).
\end{equation}
It is \emph{simple} if its only normal subgroups are the trivial
group and $G$, and \emph{finitely presented} if
$G\cong F(S)/\langle\!\langle R\rangle\!\rangle$ for finite sets
of generators $S$ and relator words $R$. We address the existence
problem asking whether these three properties can hold together
in an infinite group.

\begin{theorem}\label{thm:main}
There exists an infinite finitely presented simple amenable group.
\end{theorem}

The group will be an alternating subgroup of a polygon exchange
group on a sufficiently large finite number of squares. Both the
number of squares and the eventual sets of generators and relators
are finite. No uniform bound on their sizes is needed for the
existence statement.

\subsection{Context and the role of finite presentation}

A topological full group records the finite local orbit rules of a
dynamical system.  For a homeomorphism $\varphi$ of a Cantor space,
it consists of the homeomorphisms that agree with integer powers of
$\varphi$ on the members of a finite clopen partition.  Giordano,
Putnam and Skau developed these groups as invariants of Cantor
minimal systems and their orbit equivalence relations
\cite{GiordanoPutnamSkau1999}.  Matui proved simplicity of their
derived groups and characterized finite generation by the minimal
subshift condition \cite[Theorems~4.9 and~5.4]{Matui2006}.
Juschenko and Monod then proved amenability of the full group of
every minimal homeomorphism of a Cantor space
\cite[Theorem~A]{JuschenkoMonod2013}.  Combined with this earlier
structure theory, their Corollary~B gives infinite finitely generated
simple amenable groups.

Finite presentation is the obstruction left by those examples.
The existence question predates them: Button records it in
\cite[p.~729]{Button2010}, and Juschenko and Monod raise it again
on page~776 of \cite{JuschenkoMonod2013}.
Matui proved that the derived groups arising from minimal subshifts
are not finitely presented \cite[Theorem~5.7]{Matui2006}.
Grigorchuk and Medynets established a broader approximation property:
the full group of every Cantor minimal system is locally embeddable
into finite groups \cite[Theorem~2.6]{GrigorchukMedynets2014}.
This means that each finite subset admits an injective map into a
finite group preserving all products that stay within that subset.
The property passes to subgroups, and a finitely presented group
with this property is residually finite.  An infinite simple group
cannot be residually finite, since a nontrivial homomorphism from
it to a finite group would be injective.  Thus this local finite
approximation property also explains why the classical full-group
examples cannot settle the finite-presentation question.

There is also a presentation theory for the classical examples.
Grigorchuk and Medynets describe their derived groups using local
$3$-cycles and relations expressing refinement and disjointness of
clopen pieces \cite{GrigorchukMedynets2018}. For a minimal subshift,
their presentation can use a finite generating set, but its relations
remain infinitely indexed. Such presentations separate the task of
finding finite generators from the harder task of deriving all local
permutation laws from finitely many relations.

Subsequent constructions have strengthened the finitely generated
existence theorem in other directions. Nekrashevych constructed
infinite finitely generated simple torsion groups of intermediate
growth \cite[Theorem~1.2]{Nekrashevych2018}, and Kionke and Schesler
obtained two-generated infinite simple amenable groups
\cite{KionkeSchesler2024}. These advances do not give finite
presentation. The latter requirement still appears in Zaremsky's
July 2026 problem list \cite[Problem~10]{Zaremsky2026};
Theorem~\ref{thm:main} answers the existence question positively.

Higher-dimensional piecewise translation groups offer a broader
setting, but bring a separate amenability problem.  Chornyi,
Juschenko and Nekrashevych extended finite-generation methods to
the derived subgroups of full groups of minimal faithful $\ZZ^d$-actions
conjugate to finite-alphabet subshifts
\cite[Theorem~1 in the author version]{ChornyiJuschenkoNekrashevych2020}.
Nekrashevych's alternating full groups organize the local even
permutations through multisections, and their simplicity follows
from minimality for effective {\'e}tale groupoids with Cantor unit space
\cite[Section~3 and Theorem~4.1]{Nekrashevych2019}.
Our local even permutations are an instance of that construction; we
prove the needed comparison and simplicity statements directly.
Amenability does not follow from the rank of the acting group:
Elek and Monod constructed a free minimal Cantor $\ZZ^2$-action
whose full group contains a nonabelian free group
\cite[Theorem~1]{ElekMonod2013}.

The geometric construction has a particularly close predecessor
in the Penrose model of Chornyi, Juschenko and Nekrashevych
\cite[Section~4 in the author version]{ChornyiJuschenkoNekrashevych2020}.
They split the plane along five arithmetic families of lines, realize
polygonal pieces as clopen sets, and describe the full group by
piecewise translations on four pentagonal windows. The translation
group has rank four over $\ZZ$. Our model likewise combines a split
plane with finitely many polygonal regions, but uses four directions
over a real quadratic ring. Amenability requires a separate argument:
the Penrose example in that source is not asserted to be amenable.

Cornulier and Lacourte study the unrestricted rectangle exchange
group on $[0,1)^d$, for every $d\geq1$, allowing arbitrary real
endpoints and translation vectors \cite{CornulierLacourte2025}.
They prove that its derived subgroup is simple and ask about
amenability, second homology, and bounded relator length over the
infinite generating set of all restricted shuffles. Our group
restricts the arithmetic parameters and permits two inclined edge
directions in addition to the coordinate directions. Finite unions
of coordinate rectangles do not give these inclined boundaries.
Moreover, bounded relator length over an infinite set of generators
is a different requirement from the finite presentation sought here.
The proof below propagates relations while keeping the generating
set fixed and finite.

The amenability argument is related to the almost-invariant finite
configurations of Juschenko and Monod
\cite[Theorem~C]{JuschenkoMonod2013} and to the extensive-amenability
methods of Juschenko, Matte Bon, Monod and de la Salle
\cite[Section~5]{JuschenkoMatteBonMonodDeLaSalle2018}.
For interval-exchange groups whose translation increments modulo
one generate an abelian group of rational rank at most two,
recurrence gives an extensively amenable action on the circle.
Comparing the two continuity conventions at discontinuities gives
a cocycle into finitely supported permutations with amenable
kernel. Polygonal boundaries
are line segments, so this finite-discontinuity argument does not
apply directly. Here correlated random fields produce polygonal
partitions whose barriers respect every prescribed translation
chart. Matching cells of the same shape then yields almost-invariant
measures on the exchange group.

The homological argument follows the stability and symmetric
monoidal category approach used by Szymik and Wahl for
Higman--Thompson groups \cite{SzymikWahl2019}. Li's general
framework also reaches individual full groups and their derived
groups: for minimal ample groupoids with locally compact Hausdorff
unit space, no isolated units, and comparison, their homology is
identified with that of an infinite loop space and its universal
cover, respectively \cite[Theorem~5.18 and Corollary~5.20]{Li2025}.
We give a concrete low-degree implementation for the polygon
algebra. An explicit resolution reduces the calculation to
translation modules of polygonal step functions, and a fixed
finite collection of square copies controls the passage to the alternating
group's multiplier. Polyhedral indicator modules and finite
translation-orbit data also underlie the projection-pattern
calculations recalled in
\cite[Section~3.1]{GahlerHuntonKellendonk2013}; our coefficient
finiteness is proved directly by corner data. We use
Segal's bar and realization results and
ordinary integral group completion, with the corrected proof of
Miller and Palmer
\cite{Segal1974,McDuffSegal1976,MillerPalmer2015}.

\subsection{The construction and the proof}

The polygons have edges on level lines of
\[
 x,\quad y,\quad y-\lambda x,\quad x-\lambda y,
\]
at levels in $\OO=\ZZ[\tau]$, where $\tau=(1+\sqrt5)/2$ and
$\lambda$ is a sufficiently large power of $\tau$. The other real
embedding makes $\lambda$ small. We identify polygonal sets that
differ only on their edges, or equivalently use the compact Stone
space $X$ of their Boolean algebra on a square. For a finite
integer $m$, the group $F_m$ consists of all finite piecewise
translations of $m$ disjoint copies of $X$. Its subgroup $A_m$ is
generated by even permutations of disjoint translated polygonal
pieces, with the pieces identified by their translations.
These definitions are made precise in Section~\ref{sec:geometry}.
We eventually choose one finite $m$ satisfying all homology and
local-permutation requirements, and use that same $A_m$ for every
property.  The homological stabilization below is a tool for studying
this fixed group; the example is not a union over increasing track
numbers.

Strict area comparison supplies spare pieces. It implies that
$A_m=[F_m,F_m]$ is infinite, perfect, and simple
(Theorem~\ref{thm:simplicity}). The remaining proof has three parts.
\begin{enumerate}[leftmargin=*,label=\textup{\arabic*.}]
\item \emph{Amenability.}
For a fixed finite collection of exchanges, we build partition
boundaries from marked segments on allowable lines. One random
field selects subdivision points on each line, and a second
selects the segments between them. The fields are correlated over
a range of scales in the conjugate embedding. The correlations make
bounded translations inexpensive in total variation while keeping
individual fluctuations small. Prescribed mean values force the
chart-edge intersections to be marked and the chart boundaries
to be barriers on an event of arbitrarily high probability.
On that event every partition cell stays in a translation chart,
so its image has the same shape. Uniform matchings between cells
of the same shape convert the partition laws into finitely
supported almost invariant group measures, and then into the
F\o lner sets in
\eqref{eq:folner-introduction}
(Theorem~\ref{thm:amenability}).

\item \emph{Finite generation of the multiplier.}
We prove that the Schur multiplier $H_2(A_m;\ZZ)$ is finitely
generated once $m$ is sufficiently large. An ordered-embedding complex gives homological
stability for $F_m$ in the needed degrees. A symmetric monoidal
category of polygons then reduces its stable homology to translation
homology of integer-valued polygonal functions. Such functions are
detected by finitely many translation types of corner data, giving
the required algebraic finiteness. An argument using a fixed finite
bank of reserved tracks transfers the result to $A_m$ (Theorem~\ref{thm:multiplier}).

\item \emph{A finitely presented central cover.}
We use permutations of the tracks applied only over specified
polygonal sets, together with translations whose total shift over
the tracks is zero. Finitely many such elements generate $A_m$.
A finite collection of their true relations presents a group
$\widehat G$ mapping onto $A_m$; the task is to prove that the
kernel is central.

Start with the permutation relations for coordinate cuts at
$i\tau$ modulo one, with $i$ ranging over a fixed finite interval
of integers. To extend this
range, the large ordinary slope of $\lambda$ places an inclined
line between two small coordinate rectangles. Its small conjugate
value keeps the new transverse coordinate indices within ranges covered
by the preceding induction step. Fixed overlap relations compare
the permutation actions along that line, proving that permutations
on the two rectangles commute. This enlarges the coordinate laws.
The resulting rectangular decompositions then reduce arbitrary
polygonal tests to finitely many configurations of concurrent lines.
Finally, exact conjugation formulas for permutations of disjoint
translated pieces show that a word acting trivially on the $m$ squares
commutes with every
generator of $\widehat G$ (Theorem~\ref{thm:central-cover}).
\end{enumerate}

The last two parts have different purposes. A finitely generated
second homology group alone does not imply finite presentation.
Together with the central extension
\[
 1\longrightarrow K\longrightarrow\widehat G
 \longrightarrow A_m\longrightarrow1,
\]
the homological five-term sequence shows that $K$ is a finitely
generated abelian group. Killing finitely many generators of $K$
then gives a finite presentation of $A_m$.

Sections~\ref{sec:geometry}--\ref{sec:simplicity} construct the group
and establish simplicity. Sections~\ref{sec:amenability}
and~\ref{sec:homology} prove the first two main ingredients.
Sections~\ref{sec:lifting}--\ref{sec:transport} construct the central
cover, and Section~\ref{sec:conclusion} completes the proof of
Theorem~\ref{thm:main}.
